% =====================================================================
\chapter{Funciones inyectivas, sobreyectivas y biyectivas}
\label{tema:3}
% =====================================================================

Este tema es el núcleo práctico del módulo: decidir si una función es inyectiva, sobreyectiva o biyectiva y \emph{demostrarlo}. Sigue la sección~12.2 de Hammack~\cite{hammack}, la sección~5.2 de Velleman~\cite{velleman} y las secciones~4.3 y~4.5 de~\cite{mcs}. Los ejercicios de clasificación siguen el problema~1 de la hoja~3 de Stanford CS103~\cite{cs103ps3}.

\section{Definiciones y método de demostración}

\begin{definicion}\label{def:t3-iny}
Sea $f\colon A\to B$.
\begin{enumerate}
  \item $f$ es \emph{inyectiva} si elementos distintos tienen imágenes distintas: para todos $a,a'\in A$, si $f(a)=f(a')$ entonces $a=a'$.
  \item $f$ es \emph{sobreyectiva} si todo elemento del codominio es imagen de alguno: para todo $b\in B$ existe $a\in A$ con $f(a)=b$; equivalentemente, $f(A)=B$.
  \item $f$ es \emph{biyectiva} si es inyectiva y sobreyectiva.
\end{enumerate}
\end{definicion}

En términos de fibras: $f$ es inyectiva si cada fibra $f^{-1}(\{b\})$ tiene \emph{a lo sumo} un elemento, sobreyectiva si tiene \emph{al menos} uno y biyectiva si tiene \emph{exactamente} uno. La \cref{fig:t3-flechas} muestra los cuatro casos posibles con diagramas de flechas.

\begin{figure}[h]
\centering
\begin{tikzpicture}[>={Stealth[length=1.6mm]},font=\small,
  pt/.style={circle,fill,inner sep=1.2pt}]
  \foreach \col/\tit/\na/\nb/\arr in {
    0/{inyectiva,\\ no sobreyectiva}/2/3/{1/1,2/3},
    3.4/{sobreyectiva,\\ no inyectiva}/3/2/{1/1,2/1,3/2},
    6.8/{biyectiva}/3/3/{1/2,2/3,3/1},
    10.2/{ninguna}/3/3/{1/1,2/1,3/3}} {
    \begin{scope}[xshift=\col cm]
      \foreach \i in {1,...,\na} \node[pt] (a\i) at (0,-0.6*\i) {};
      \foreach \j in {1,...,\nb} \node[pt] (b\j) at (1.4,-0.6*\j) {};
      \draw[rounded corners] (-0.3,-0.25) rectangle (0.3,-0.6*\na-0.35);
      \draw[rounded corners] (1.1,-0.25) rectangle (1.7,-0.6*\nb-0.35);
      \foreach \u/\v in \arr \draw[->] (a\u) -- (b\v);
      \node[align=center,anchor=north] at (0.7,-2.3) {\tit};
    \end{scope}}
\end{tikzpicture}
\caption{Diagramas de flechas de funciones entre conjuntos finitos.}
\label{fig:t3-flechas}
\end{figure}

\paragraph{Cómo demostrar cada propiedad.}
\begin{itemize}[leftmargin=1.4em]
  \item \emph{Inyectividad}: se supone $f(a)=f(a')$ para $a,a'$ arbitrarios y se deduce $a=a'$ (demostración directa), o se supone $a\neq a'$ y se deduce $f(a)\neq f(a')$ (contrarrecíproco).
  \item \emph{No inyectividad}: basta un contraejemplo, $a\neq a'$ con $f(a)=f(a')$.
  \item \emph{Sobreyectividad}: se toma $b\in B$ arbitrario, se \emph{busca} un candidato $a$ (normalmente resolviendo la ecuación $f(a)=b$ en borrador), y se \emph{comprueba} que $a\in A$ y que $f(a)=b$.
  \item \emph{No sobreyectividad}: basta un $b\in B$ y una prueba de que la ecuación $f(a)=b$ no tiene solución en $A$.
\end{itemize}

\begin{ejemplo}\label{ej:t3-2n1}
Sea $f\colon\Z\to\Z$, $f(n)=2n+1$. \emph{Es inyectiva}: si $2n+1=2m+1$, restando~1 y dividiendo entre~2 se obtiene $n=m$. \emph{No es sobreyectiva}: $f(n)$ es siempre impar, luego $2$ no tiene preimagen. En cambio, $g\colon\R\to\R$, $g(x)=2x+1$, es biyectiva: la inyectividad se prueba igual y, dado $y\in\R$, el número $x=(y-1)/2$ es real y cumple $g(x)=y$. El mismo cálculo falla en $\Z$ porque $(y-1)/2$ no es entero cuando $y$ es par: la sobreyectividad depende del dominio.
\end{ejemplo}

\begin{ejemplo}[Dominio y codominio importan]\label{ej:t3-cuadrado}
La regla $x\mapsto x^{2}$ da cuatro funciones distintas:
\begin{center}\small
\begin{tabular}{@{}lccl@{}}
\toprule
Función & Inyectiva & Sobreyectiva & Motivo \\
\midrule
$\R\to\R$ & no & no & $f(-1)=f(1)$; $-1$ no tiene preimagen \\
$\R\to[0,\infty)$ & no & sí & $f(-1)=f(1)$; $y\ge0$ es $f(\sqrt y)$ \\
$[0,\infty)\to\R$ & sí & no & $x^{2}=x'^{2}$ con $x,x'\ge0$ da $x=x'$ \\
$[0,\infty)\to[0,\infty)$ & sí & sí & las dos razones anteriores \\
\bottomrule
\end{tabular}
\end{center}
\end{ejemplo}

\begin{ejemplo}[Una función de dos variables]\label{ej:t3-suma}
Sea $s\colon\Z\times\Z\to\Z$, $s(m,n)=m+n$. Es sobreyectiva, pues $s(k,0)=k$ para todo $k\in\Z$, y no es inyectiva, pues $s(1,0)=s(0,1)$ con $(1,0)\neq(0,1)$.
\end{ejemplo}

\begin{ejemplo}[Funciones sobre cadenas]\label{ej:t3-cadenas}
Sea $\Sigma$ un alfabeto. La función $c\colon\Sigma^{*}\to\Sigma^{*}$, $c(w)=\cad{a}w$ (anteponer el símbolo $\cad{a}\in\Sigma$), es inyectiva: si $\cad{a}w=\cad{a}w'$, al suprimir el primer símbolo de ambas cadenas queda $w=w'$. No es sobreyectiva: la cadena vacía $\vacia$ no empieza por $\cad{a}$. La longitud $\ell\colon\Sigma^{*}\to\N$ es sobreyectiva si $\Sigma\neq\emptyset$ (la cadena $\cad{a}^{n}$ tiene longitud $n$) y no es inyectiva si $\abs{\Sigma}\ge2$.
\end{ejemplo}

\section{Criterios para funciones reales}

\begin{definicion}
$f\colon A\to\R$, con $A\subseteq\R$, es \emph{estrictamente creciente} si $x<x'$ implica $f(x)<f(x')$, y \emph{estrictamente decreciente} si $x<x'$ implica $f(x)>f(x')$. Es \emph{estrictamente monótona} si es una de las dos.
\end{definicion}

\begin{proposicion}\label{prop:t3-monotona}
Toda función estrictamente monótona es inyectiva.
\end{proposicion}

\begin{proof}
Sea $f$ estrictamente creciente y sean $x\neq x'$. Sin pérdida de generalidad, $x<x'$; entonces $f(x)<f(x')$, luego $f(x)\neq f(x')$. El caso decreciente es igual.
\end{proof}

Si $f$ es derivable en un intervalo con $f'>0$ en todo punto, entonces es estrictamente creciente (consecuencia del teorema del valor medio, que se estudia en el módulo de cálculo de esta fase), y por tanto inyectiva. Para la sobreyectividad sobre un intervalo, el teorema de Bolzano (valores intermedios) permite concluir que una función continua toma todos los valores comprendidos entre sus límites; en este módulo se prefiere resolver $f(x)=y$ explícitamente.

\begin{ejemplo}[La sigmoide]\label{ej:t3-sigmoide}
Sea $\sigma\colon\R\to(0,1)$, $\sigma(x)=1/(1+e^{-x})$. Es inyectiva: si $\sigma(x)=\sigma(x')$, entonces $1+e^{-x}=1+e^{-x'}$, luego $e^{-x}=e^{-x'}$ y, por la inyectividad de la exponencial, $x=x'$. Es sobreyectiva sobre $(0,1)$: dado $y\in(0,1)$, se resuelve $1+e^{-x}=1/y$, es decir, $e^{-x}=(1-y)/y>0$, y $x=\ln\bigl(y/(1-y)\bigr)$ es un número real que cumple $\sigma(x)=y$. La sigmoide es la función de activación del módulo~11 y su inversa es el \emph{logit}.
\end{ejemplo}

\section{Funciones entre conjuntos finitos}

\begin{proposicion}\label{prop:t3-finitos}
Sean $A$ y $B$ finitos y $f\colon A\to B$.
\begin{enumerate}
  \item Si $f$ es inyectiva, $\abs{A}\le\abs{B}$; si es sobreyectiva, $\abs{A}\ge\abs{B}$.
  \item Si $\abs{A}=\abs{B}$, entonces $f$ es inyectiva si y solo si es sobreyectiva.
\end{enumerate}
\end{proposicion}

\begin{proof}
(i) El dominio es la unión disjunta de las fibras: $\abs{A}=\sum_{b\in B}\abs{f^{-1}(\{b\})}$. Si $f$ es inyectiva, cada sumando es $\le1$ y la suma es $\le\abs{B}$; si es sobreyectiva, cada sumando es $\ge1$ y la suma es $\ge\abs{B}$. (ii) Si $f$ es inyectiva, las fibras tienen a lo sumo un elemento y $\abs{f(A)}=\abs{A}=\abs{B}$; como $f(A)\subseteq B$ y ambos son finitos con el mismo cardinal, $f(A)=B$. El recíproco es el \cref{eje:3-finitos}.
\end{proof}

\begin{proposicion}\label{prop:t3-contar-iny}
Si $\abs{A}=k$ y $\abs{B}=n$, el número de funciones inyectivas $A\to B$ es $n(n-1)\cdots(n-k+1)$, que vale~0 si $k>n$. En particular, hay $n!$ biyecciones entre dos conjuntos de $n$ elementos.
\end{proposicion}

\begin{proof}
Sea $A=\{a_1,\dots,a_k\}$. Para $f(a_1)$ hay $n$ elecciones; para $f(a_2)$, $n-1$ (debe ser distinto del anterior); y así hasta $f(a_k)$, con $n-k+1$ elecciones. Si $k>n$ alguno de los factores es cero.
\end{proof}

\begin{observacion}[Conexión con el curso]
Un tokenizador debe ser \emph{inyectivo}: dos textos distintos no pueden producir la misma secuencia de tokens, o el texto original no podría recuperarse (módulo~1, §1.2). La codificación UTF-8 es inyectiva y ningún código es prefijo de otro, lo que garantiza que una cadena de bytes se descompone de una sola manera (módulo~12, §1.4). La asignación de índices a los tokens, $\V\to\{0,\dots,K-1\}$, es una biyección. En cambio, la decodificación de tokens a texto \emph{no} es inyectiva en general: las secuencias $(\cad{a},\cad{b})$ y $(\cad{ab})$ producen el mismo texto (\cref{eje:3-codigos}).
\end{observacion}

\begin{observacion}[Errores frecuentes]
(1) «Demostrar» la inyectividad partiendo de $a=a'$ y llegando a $f(a)=f(a')$: eso prueba solo que $f$ está bien definida. (2) Comprobar unos cuantos valores en lugar de razonar con elementos arbitrarios. (3) En la sobreyectividad, no comprobar que el candidato pertenece al dominio: $(y-1)/2$ resuelve $2n+1=y$ pero no es entero si $y$ es par. (4) Olvidar que inyectividad y sobreyectividad dependen del dominio y del codominio, no solo de la fórmula (\cref{ej:t3-cuadrado}).
\end{observacion}

\begin{lecturas}
\item Hammack~\cite{hammack}, §12.2 «Injective and Surjective Functions».
\item Velleman~\cite{velleman}, §5.2 «One-to-One and Onto».
\item Lehman, Leighton y Meyer~\cite{mcs}, §4.3 «Functions» y §4.5 «Finite Cardinality».
\item Stanford CS103~\cite{cs103ps3}, Problem Set~3, problemas 1 a~3.
\end{lecturas}

\section*{Ejercicios del tema 3}
\addcontentsline{toc}{section}{Ejercicios del tema 3}

\begin{ejercicio}{1}{eje:3-tablas}
Clasifique (inyectiva, sobreyectiva, biyectiva o ninguna) las funciones $\{1,2,3\}\to\{\cad{a},\cad{b},\cad{c}\}$ dadas por: (a) $1\mapsto\cad{a}$, $2\mapsto\cad{b}$, $3\mapsto\cad{c}$; (b) $1\mapsto\cad{a}$, $2\mapsto\cad{a}$, $3\mapsto\cad{c}$; (c) $1\mapsto\cad{c}$, $2\mapsto\cad{a}$, $3\mapsto\cad{b}$; (d) $1\mapsto\cad{b}$, $2\mapsto\cad{b}$, $3\mapsto\cad{b}$. Para las que no sean biyectivas, ¿podría ser alguna inyectiva sin ser sobreyectiva?
\end{ejercicio}

\begin{ejercicio}{1}{eje:3-error}
Se propone la siguiente «demostración» de que $f\colon\Z\to\Z$, $f(n)=n^{2}$, es inyectiva: «Sean $a,b\in\Z$ con $a=b$. Entonces $a^{2}=b^{2}$, es decir, $f(a)=f(b)$. Luego $f$ es inyectiva.» Señale el error y decida si $f$ es inyectiva.
\end{ejercicio}

\begin{ejercicio}{2}{eje:3-clasificar}
Clasifique, justificando brevemente: (a) $\Z\to\Z$, $n\mapsto n+3$; (b) $\Z\to\Z$, $n\mapsto 2n$; (c) $\Z\to\Z$, $n\mapsto n^{3}$; (d) $\N\to\N$, $n\mapsto n+1$; (e) $\R\to\R$, $x\mapsto x^{3}-x$; (f) $\R\to\R$, $x\mapsto e^{x}$; (g) $\R\to\Z$, $x\mapsto\lfloor x\rfloor$; (h) $\Z\to\N$, $n\mapsto\abs{n}$.
\fuente{adaptado de Stanford CS103~\cite{cs103ps3}, Problem Set~3, problema~1 «Properties of Functions».}
\end{ejercicio}

\begin{ejercicio}{2}{eje:3-5n2}
Sea $f\colon\Z\to\Z$, $f(n)=5n+2$. Demuestre que $f$ es inyectiva y no sobreyectiva. ¿Qué ocurre con la misma fórmula como función $\Q\to\Q$?
\fuente{adaptado de Hammack~\cite{hammack}, §12.2, ejercicios de la sección.}
\end{ejercicio}

\begin{ejercicio}{2}{eje:3-homografica}
Sea $f\colon\R\setminus\{2\}\to\R\setminus\{3\}$, $f(x)=\dfrac{3x+1}{x-2}$. (a) Compruebe que $f(x)\neq3$ para todo $x$ del dominio, de modo que el codominio es correcto. (b) Demuestre que $f$ es biyectiva y obtenga la expresión de la única preimagen de cada $y$.
\fuente{adaptado de University of Washington, Math 300~\cite{uw300}, «Problems on functions», problema~5.}
\end{ejercicio}

\begin{ejercicio}{3}{eje:3-lineal-z2}
Sea $f\colon\Z\times\Z\to\Z\times\Z$, $f(m,n)=(m+n,\ m+2n)$. Demuestre que $f$ es biyectiva.
\fuente{adaptado de Hammack~\cite{hammack}, §12.2, ejercicios de la sección.}
\end{ejercicio}

\begin{ejercicio}{3}{eje:3-2m3n}
Sea $f\colon\Z\times\Z\to\Z$, $f(m,n)=2m+3n$. Demuestre que $f$ es sobreyectiva y no inyectiva. ¿Es sobreyectiva $g(m,n)=2m+4n$?
\fuente{adaptado de Hammack~\cite{hammack}, §12.2, ejercicios de la sección.}
\end{ejercicio}

\begin{ejercicio}{3}{eje:3-NZ}
Sea $f\colon\N\to\Z$ definida por $f(n)=n/2$ si $n$ es par y $f(n)=-(n+1)/2$ si $n$ es impar. (a) Calcule $f(0),\dots,f(6)$. (b) Demuestre que $f$ es biyectiva.
\fuente{adaptado de Stanford CS103~\cite{cs103ps3}, Problem Set~3, problema~2.}
\end{ejercicio}

\begin{ejercicio}{3}{eje:3-creciente}
(a) Decida si $x\mapsto x^{2}$ y $x\mapsto x^{3}$ son estrictamente crecientes en $\R$. (b) Demuestre que si $f,g\colon\R\to\R$ son estrictamente crecientes, también lo es $g\circ f$ (la composición se define en el \cref{tema:4}: $(g\circ f)(x)=g(f(x))$). (c) ¿Es cierto que toda función inyectiva $\R\to\R$ es estrictamente monótona? Demuéstrelo o dé un contraejemplo.
\fuente{adaptado de Stanford CS103~\cite{cs103ps3}, Problem Set~3, problema~3 «Strictly Increasing Functions».}
\end{ejercicio}

\begin{ejercicio}{3}{eje:3-cadenas}
Sea $\Sigma=\{\cad{a},\cad{b}\}$. Para cada función $\Sigma^{*}\to\Sigma^{*}$, decida si es inyectiva y si es sobreyectiva, y demuéstrelo: (a) la inversión $r(w_1w_2\cdots w_n)=w_n\cdots w_2w_1$; (b) la duplicación $d(w)=ww$; (c) la función $q$ que suprime el último símbolo de una cadena no vacía y deja igual la cadena vacía.
\end{ejercicio}

\begin{ejercicio}{3}{eje:3-codigos}
Sea $\Sigma$ un alfabeto finito, $k\ge1$ y $c\colon\Sigma\to\{0,1\}^{k}$ una función inyectiva (un \emph{código de longitud fija}). Su extensión $c^{*}\colon\Sigma^{*}\to\{0,1\}^{*}$ concatena los códigos: $c^{*}(x_1\cdots x_n)=c(x_1)\cdots c(x_n)$. (a) Demuestre que $c^{*}$ es inyectiva. (b) Sea ahora el código de longitud variable $c(\cad{x})=0$, $c(\cad{y})=01$, $c(\cad{z})=1$. Demuestre que $c$ es inyectiva pero $c^{*}$ no lo es. (c) Relacione (b) con la decodificación de un vocabulario que contiene \cad{a}, \cad{b} y \cad{ab}.
\end{ejercicio}

\begin{ejercicio}{2}{eje:3-contar}
Calcule el número de funciones, de funciones inyectivas, de sobreyectivas y de biyectivas: (a) de $\{1,2,3\}$ en $\{1,2,3,4\}$; (b) de $\{1,2,3,4\}$ en $\{1,2,3\}$; (c) de $\{1,2,3\}$ en sí mismo. Para el número de sobreyecciones de~(b), razone qué elemento de la imagen tiene dos preimágenes.
\end{ejercicio}

\begin{ejercicio}{3}{eje:3-finitos}
Sean $A$ y $B$ conjuntos finitos con $\abs{A}=\abs{B}$ y $f\colon A\to B$ sobreyectiva. Demuestre que $f$ es inyectiva. Muestre con un ejemplo que el resultado es falso para conjuntos infinitos.
\fuente{adaptado de Lehman, Leighton y Meyer~\cite{mcs}, §4.5 «Finite Cardinality».}
\end{ejercicio}

\begin{ejercicio}{3}{eje:3-softsign}
Sea $f\colon\R\to(-1,1)$, $f(x)=\dfrac{x}{1+\abs{x}}$. (a) Compruebe que $f(x)\in(-1,1)$ para todo $x$. (b) Demuestre que $f$ es inyectiva. (c) Demuestre que es sobreyectiva y halle la preimagen de cada $y\in(-1,1)$.
\end{ejercicio}

\begin{ejercicio}{2}{eje:3-tanh}
Sea $\tanh\colon\R\to(-1,1)$, $\tanh x=\dfrac{e^{x}-e^{-x}}{e^{x}+e^{-x}}$. Demuestre que es biyectiva resolviendo $\tanh x=y$ y compruebe que la preimagen de $y$ es $\tfrac12\ln\dfrac{1+y}{1-y}$. Calcule la preimagen de $y=0{,}6$.
\end{ejercicio}

\begin{ejercicio}{1}{eje:3-fibras}
Sea $f\colon A\to B$ con $A$ y $B$ finitos. Explique, en términos del tamaño de las fibras $f^{-1}(\{b\})$, qué significa que $f$ sea inyectiva, sobreyectiva o biyectiva, y por qué $\abs{A}=\sum_{b\in B}\abs{f^{-1}(\{b\})}$.
\end{ejercicio}
